\documentclass[a4paper,12pt]{article}\usepackage[]{graphicx}\usepackage[]{color}
%% maxwidth is the original width if it is less than linewidth
%% otherwise use linewidth (to make sure the graphics do not exceed the margin)
\makeatletter
\def\maxwidth{ %
  \ifdim\Gin@nat@width>\linewidth
    \linewidth
  \else
    \Gin@nat@width
  \fi
}
\makeatother

\definecolor{fgcolor}{rgb}{0.345, 0.345, 0.345}
\newcommand{\hlnum}[1]{\textcolor[rgb]{0.686,0.059,0.569}{#1}}%
\newcommand{\hlstr}[1]{\textcolor[rgb]{0.192,0.494,0.8}{#1}}%
\newcommand{\hlcom}[1]{\textcolor[rgb]{0.678,0.584,0.686}{\textit{#1}}}%
\newcommand{\hlopt}[1]{\textcolor[rgb]{0,0,0}{#1}}%
\newcommand{\hlstd}[1]{\textcolor[rgb]{0.345,0.345,0.345}{#1}}%
\newcommand{\hlkwa}[1]{\textcolor[rgb]{0.161,0.373,0.58}{\textbf{#1}}}%
\newcommand{\hlkwb}[1]{\textcolor[rgb]{0.69,0.353,0.396}{#1}}%
\newcommand{\hlkwc}[1]{\textcolor[rgb]{0.333,0.667,0.333}{#1}}%
\newcommand{\hlkwd}[1]{\textcolor[rgb]{0.737,0.353,0.396}{\textbf{#1}}}%
\let\hlipl\hlkwb


\usepackage{framed}
\makeatletter
\newenvironment{kframe}{%
 \def\at@end@of@kframe{}%
 \ifinner\ifhmode%
  \def\at@end@of@kframe{\end{minipage}}%
  \begin{minipage}{\columnwidth}%
 \fi\fi%
 \def\FrameCommand##1{\hskip\@totalleftmargin \hskip-\fboxsep
 \colorbox{shadecolor}{##1}\hskip-\fboxsep
     % There is no \\@totalrightmargin, so:
     \hskip-\linewidth \hskip-\@totalleftmargin \hskip\columnwidth}%
 \MakeFramed {\advance\hsize-\width
   \@totalleftmargin\z@ \linewidth\hsize
   \@setminipage}}%
 {\par\unskip\endMakeFramed%
 \at@end@of@kframe}
\makeatother

\definecolor{shadecolor}{rgb}{.97, .97, .97}
\definecolor{messagecolor}{rgb}{0, 0, 0}
\definecolor{warningcolor}{rgb}{1, 0, 1}
\definecolor{errorcolor}{rgb}{1, 0, 0}
\newenvironment{knitrout}{}{} % an empty environment to be redefined in TeX

\usepackage{alltt}
\usepackage[applemac]{inputenc}
\usepackage[OT1,T1]{fontenc}
% \usepackage{times}
\usepackage{setspace}
\usepackage{lscape}
\usepackage{amsmath,amssymb}
\usepackage{fancyhdr}
\usepackage{fancyvrb,moreverb,boxedminipage}
\usepackage{multirow}
\usepackage{float}
\usepackage{hyperref}
\usepackage{listings}
\usepackage{verbatim}
\usepackage{enumitem}
\usepackage{bbm}

\setlength{\textwidth}{160mm}
\setlength{\textheight}{230mm}
\setlength{\oddsidemargin}{-5mm}
\setlength{\topmargin}{-10mm}

\textwidth=6in
\textheight=8in
\topmargin=0.0in
\oddsidemargin=0in 
\baselineskip=18pt
\headheight=2cm
\setcounter{MaxMatrixCols}{10}

\DeclareMathOperator*{\argmax}{argmax}
\newcommand{\ns}{\!\!\!}
\newcommand{\e}{\varepsilon}
\newcommand{\be}{\beta}
\newcommand{\s}{\sigma}
\newcommand{\vv}{\vspace{.5cm}}
\newcommand{\beps}{\boldsymbol{\epsilon}}
\newcommand{\sumin}{\sum_{i=1}^n}
\newcommand{\bg}[1]{\mbox{\boldmath{$#1$}}}
\DeclareMathOperator{\plim}{plim}
\DefineVerbatimEnvironment{code}{Verbatim}{fontsize=\small}
\DefineVerbatimEnvironment{example}{Verbatim}{fontsize=\small}

\def\by{\mathbf{y}}
\def\bx{\mathbf{x}}
\def\ba{\mathbf{a}}
\def\bb{\mathbf{b}}
\def\be{\mathbf{e}}
\def\bu{\mathbf{u}}
\def\bv{\mathbf{v}}
\def\bz{\mathbf{z}}

%===========
%Greek letter vectors
%===========
\def\bbeta{\mbox{\boldmath $\beta$}}
\def\bgamma{\mbox{\boldmath $\gamma$}}
\def\bvareps{\mbox{\boldmath $\varepsilon$}}
\def\bleta{\mbox{\boldmath $\eta$}}
\def\bmu{\mbox{\boldmath $\mu$}}
\def\btheta{\mbox{\boldmath $\theta$}}
\def\bsigma{\mbox{\boldmath $\sigma$}}
\def\blambgam{\mbox{\boldmath $\lambda_{\gamma}$}}
\def\blambbet{\mbox{\boldmath $\lambda_{\beta}$}}
\def\blambda{\mbox{\boldmath $\lambda$}}

%===========
%Matrices
%===========
\def\bX{\mathbf{X}}
\def\bD{\mathbf{D}(\bsigma)}
\def\bV{\mathbf{V}(\bsigma)}
\def\bVinv{\mathbf{V}^{-1}(\bsigma)}
\def\bR{\mathbf{R}(\bsigma)}
\def\bA{\mathbf{A}}
\def\bB{\mathbf{B}}
\def\bZ{\mathbf{Z}}
\def\bIn{\mathbf{I_{N}}}
\def\bIm{\mathbf{I_{m}}}
\def\bIni{\mathbf{I_{n}}}
\def\SSW{\mathrm{SSW}}
\def\SSB{\mathrm{SSB}}
\def\Normal{\mathcal{N}\!}

%===========
%Parentheses, etc.
%===========
\def\op{\left(}
\def\cp{\right)}
\def\oc{\left[}
\def\cc{\right]}
\def\ok{\left\{}
\def\ck{\right\}}

%===========
%Statistical functions and real numbers.
%===========
\def\exE{\mathbb{E}}
\def\Real{\mathbb{R}}
\def\Var{\mathbb{V}ar}
\def\Proba{\mathbb{P}}

%===========
%Listings and hyperlink settings
%===========
\hypersetup{colorlinks=true,linkcolor=blue,citecolor=blue, urlcolor=blue}
\urlstyle{tt}
\lstset{language=R, basicstyle=\footnotesize\ttfamily,}

%===========
%Special Environments
%===========
\newenvironment{exercices}{\newcounter{moncompteur}\begin{list}
   {\bfseries Ex. \arabic{moncompteur}}{\usecounter{moncompteur}
     \setlength{\labelwidth}{2.6em}\setlength{\leftmargin}{3.1em}}}{\end{list}}
\renewcommand{\labelenumi}{\bfseries\alph{enumi})}
\newcommand{\splus}[3][0]{\texttt{>~#2}\ \ \ \ \textit{#3}\\[#1mm]}

\newenvironment{exo}
{\begin{center}\setlength{\textwidth}{140mm} \underline{Exercise:} \begin{minipage}[t]{120mm}}
{\end{minipage}\setlength{\textwidth}{160mm}\end{center}}


\usepackage{verbatim}

\lstset{ %
  language=R, 
  basicstyle=\ttfamily,
  backgroundcolor=\color{lightgray},   % choose the background color; you must add \usepackage{color} or \usepackage{xcolor}
  xleftmargin= 10 pt,
  xrightmargin= 10 pt,
  escapeinside={\%*}{*)},          % if you want to add LaTeX within your code
  frame=none,                    % adds a frame around the code
  numbers=none,                    % where to put the line-numbers; possible values are (none, left, right)
  numbersep=5pt,                   % how far the line-numbers are from the code
  numberstyle=\tiny\color{mygray}, % the style that is used for the line-numbers
  showspaces=false,                % show spaces everywhere adding particular underscores; it overrides 'showstringspaces'
}

\pagestyle{fancy}


\definecolor{mygreen}{rgb}{0,0.6,0}
\definecolor{mygray}{rgb}{0.5,0.5,0.5}
\definecolor{mymauve}{rgb}{0.58,0,0.82}
\definecolor{lightgray}{gray}{0.95}
\IfFileExists{upquote.sty}{\usepackage{upquote}}{}
\begin{document}

\lhead{University of Geneva\\ \textbf{GLAM}\\
Pr. Eva Cantoni}
\rhead{GSEM\\ Spring 2022 \\ \textbf{Homework 05}}

\begin{center}
\bigskip

\bigskip

\bigskip

\bigskip

\bigskip

\bigskip

\fbox{{\Large Hurdle and zero-inflated models - Handout}}

\bigskip

\bigskip

\bigskip
\end{center}

%%%%%%%%%%%%%%%%%%%%%%%%%%%%%%%%%%%%%%%%%%%%%%%%%%%%%%%%%%%%%%%%%%%%%%%%%%%%


\setlength{\parskip}{1ex plus 0.5ex minus 0.2ex}
\parindent=0cm
\setlength{\parskip}{1ex plus 0.5ex minus 0.2ex}
\linespread{1.1}
\sloppy
\section{Exercise 1}

A random variable $Y$ is said to follow a negative binomial distribution, denoted as $\mathcal{NB}(\mu,\vartheta)$, when its probability mass function (p.m.f.) is
\begin{eqnarray*}
\Proba(Y = y ; \vartheta, \mu) = \frac {\Gamma(y+\vartheta)}{\Gamma(\vartheta)\Gamma(y+1)}\left(\frac{\vartheta}{\mu+\vartheta}\right)^\vartheta
\left(\frac{\mu}{\mu+\vartheta}\right)^{y},
\end{eqnarray*}
where $\Gamma(u)=\int_0^\infty t^{u-1}\exp(-t)dt$ denotes the gamma function, $\mu>0$, $\vartheta>0$ and $y~=~0,1,2,\ldots$.
\begin{enumerate}[ref=\bf{\alph*)}]
\item Give the p.m.f. of a random variable distributed as a truncated $\mathcal{NB}(\mu,\vartheta)$, \textit{i.e.} with strictly positive outcomes.

\color{blue}
We are interested in finding the expression of :

$$
\Proba \op Y = y  |   Y> 0 ; \mu , \vartheta \cp  
= \frac {\Proba \op Y = y , y > 0 ; \mu , \vartheta \cp}{\Proba\op Y>0 ; \mu, \vartheta\cp}
= \frac {\Proba \op Y = y , y > 0 ; \mu , \vartheta \cp}{1-\Proba\op Y = 0 ; \mu, \vartheta\cp}
$$  	

It can be shown that: 

$$
\Proba\op Y=0 ; \mu, \vartheta\cp = \op \frac{\vartheta }{\mu + \vartheta }\cp^\vartheta \\ 
$$  

Yielding : 
$$
\Proba \op Y = y | Y> 0 ; \mu , \vartheta \cp 
= \frac {\Gamma(y+\vartheta)}{\Gamma(\vartheta)\Gamma(y+1)} 
\frac{\vartheta^\vartheta}{\op\mu+\vartheta\cp^\vartheta - \vartheta^\vartheta}
\op\frac{\mu}{\mu+\vartheta}\cp^{y}
$$ 

\color{black}
\item \label{loglkhd}Define a hurdle model for responses issued from a $\mathcal{NB}$ distribution with excess of zeros. Write the log-likelihood function of that model.
\color{blue} 

By defining $\pi_i = \pi\op \bx_i ; \bbeta\cp$ the probability of ``crossing the hurdle'' i.e. obtaining an observation different from zero, and  $\mu_i = \mu\op \bz_i; \bgamma\cp$ the conditional expectation of the count process, we can provide the expression for the likelihood contributions using the p.m.f. we've found on the previous point. 

$$
\Proba\op Y_i = y_i | \bx_i , \bz_i ; \bbeta, \bgamma \cp 
= 
\ok
\begin{array}{c c}
1- \pi_i  & \textrm{if : } y_i = 0 \\ 
\pi_i
\frac {\Gamma(y_i+\vartheta)}{\Gamma(\vartheta)\Gamma(y_i+1)} 
\frac{\vartheta^\vartheta}{\op\mu_i+\vartheta\cp^\vartheta - \vartheta^\vartheta}
\op\frac{\mu_i}{\mu_i+\vartheta}\cp^{y_i} & \textrm{if : } y_i > 0 
\end{array}
\right.
$$ 

Yielding the likelihood : 
$$
\mathcal{L}\op \bbeta, \bgamma  ; \by ,  \bX , \bZ  \cp 
= \prod_{i=1}^n 
\oc 1- \pi_i \cc^{\mathbbm{1}\op y_i =  0 \cp}
\oc 
\pi_i
\frac {\Gamma(y_i+\vartheta)}{\Gamma(\vartheta)\Gamma(y_i+1)} 
\frac{\vartheta^\vartheta}{\op\mu_i+\vartheta\cp^\vartheta - \vartheta^\vartheta}
\op\frac{\mu_i}{\mu_i+\vartheta}\cp^{y_i} 
\cc^{\mathbbm{1}\op y_i >  0 \cp}
$$ 
And the log-likelihood : 
\begin{align*}
\ell\op \bbeta, \bgamma  ; \by ,  \bX , \bZ  \cp 
&= \sum_{i=1}^n 
{\mathbbm{1}\op y_i =  0 \cp} \log \op 1- \pi_i \cp + {\mathbbm{1}\op y_i >  0 \cp} \log \pi_i \\
&+ 
{\mathbbm{1}\op y_i >  0 \cp}\log 
\oc 
\frac {\Gamma(y_i+\vartheta)}{\Gamma(\vartheta)\Gamma(y_i+1)} 
\frac{\vartheta^\vartheta}{\op\mu_i+\vartheta\cp^\vartheta - \vartheta^\vartheta}
\op\frac{\mu_i}{\mu_i+\vartheta}\cp^{y_i} 
\cc 
\end{align*} 

\color{black}

\item{Show that, in a similar fashion as for the Poisson hurdle model, the log-likelihood function computed in \ref{loglkhd} can be separated into two components: the first one for the peak at 0 and the second one for the strictly positive outcomes.

\color{blue} 

It is clear that we can separate terms according to their dependence on $\bbeta$ and $\bgamma$ in the last expression, yielding:

\begin{align*}
\ell \op \bbeta, \bgamma  ; \by ,  \bX , \bZ  \cp 
&= \sum_{i=1}^n 
{\mathbbm{1}\op y_i =  0 \cp} \log \op 1- \pi_i \cp + {\mathbbm{1}\op y_i >  0 \cp} \log \pi_i \\
&+ 
\sum_{i=1}^n 
{\mathbbm{1}\op y_i >  0 \cp}\log 
\oc 
\frac {\Gamma(y_i+\vartheta)}{\Gamma(\vartheta)\Gamma(y_i+1)} 
\frac{\vartheta^\vartheta}{\op\mu_i+\vartheta\cp^\vartheta - \vartheta^\vartheta}
\op\frac{\mu_i}{\mu_i+\vartheta}\cp^{y_i} 
\cc \\
\ell \op \bbeta, \bgamma  ; \by ,  \bX , \bZ  \cp 
&= \ell_1\op\bbeta ; \by, \bX\cp + \ell_2\op \bgamma, \vartheta ; \by, \bZ \cp 
\end{align*}  
\color{black}
\begin{enumerate}
\item What does it imply for the estimation process?

\color{blue} 
The estimation of parameters $\bbeta$ and $\bgamma$ can be performed separately using either of these log-likelihood components. It is customary to use a logit link to model the probability of observing a positive count and a log-link to model the count process.  

\color{black}
\item What does it imply from the point of view of statistical inference?

\color{blue} 
Both sets of estimates are uncorrelated i.e. $\mathbb{C}ov\op\widehat{\bbeta} , \widehat{\bgamma}\cp = \mathbf{0}$  \color{black}
\end{enumerate}}
\end{enumerate}

\section{Exercise 2}

Assume that $Y_1, \ldots, Y_n$ are responses for $n$ individuals issued from 
a
zero-inflated Poisson model (ZIP) with covariates $x_1, \ldots, x_n$ and
parameter $\alpha$ explaining zero responses; and covariates $z_1, \ldots,
z_n$ and parameter $\delta$ explaining positive responses.


Let $W_i$ be an unobservable binary variable that indicates the risk state of 
individual
$i$: $W_i=1$ if the subject is not at risk of an event; and $W_i=0$
if the subject is at risk. Notice that $Y_i$ denote the event count for a
subject and this event occurs only when $W_i=0$. We also suppose  that
$\Pr\{W_i=1\}=\pi_i$ (where $\pi_i$ is a function of $x_i^T\alpha$) and
$P\{Y_i=y|W_i=0\}=\lambda_i^y e^{-\lambda_i}/y!$ (where $\lambda_i$ is a
function of $z_i^T\delta$) i.e.  $Y_i|W_i=0$ is issued from the Poisson
distribution.


\begin{enumerate}
  \item[a)] Write the likelihood function $$L(\alpha, \delta |y_1, \ldots, y_n,
w_1, \ldots, w_n)=\prod_{i=1}^n P\{Y_i=y_i|W_i=w_i\}.$$ 
Show that the
corresponding log-likelihood function $\ell(\alpha, \delta |y_1, \ldots,
y_n, w_1, \ldots, w_n)$ can be written as $$\ell(\alpha |y_1, \ldots, y_n,
w_1, \ldots, w_n)+\ell(\delta |y_1, \ldots, y_n, w_1, \ldots, w_n)$$ i.e.\ the
parameters explaining the presence of zeros can be separated from the other
parameters. What does it imply numerically?\\
\end{enumerate}  
\color{blue}
\begin{align*}
    L(\alpha, \delta | y_1,...,y_n,w_1,...,w_n) &= \prod_{i=1}^n P(Y_i=y_i|W_i=w_i)\\
    &= \prod_{i=1}^n \pi(x_i^T\alpha)^{w_i} (1-\pi(x_i^T\alpha))^{1-w_i}\left(\frac{\exp(-\lambda (z_i^T\delta))\lambda(z_i^T\delta)^{y_i}}{y_i!}\right)^{1-w_i}
\end{align*}
Then
\begin{align*}
    l(\alpha, \delta | y_1,...,y_n,w_1,...,w_n)=&\sum_{i=1}^n \{w_i\log\pi(x_i^T\alpha) + (1-w_i)\log(1-\pi(x_i^T\alpha))\} + \\
    &\sum_{i=1}^n\{-(1-w_i)\lambda(z_i^T\delta)+(1-w_i)y_i\log\lambda(z_i^T\delta)\}-\\
    &\sum_{i=1}^n(1-w_i)\log y_i!\\
    =&l(\alpha | y_1,...,y_n,w_1,...,w_n) + l(\delta | y_1,...,y_n,w_1,...,w_n) - \\
    &\sum_{i=1}^n (1-w_i)\log y_i!
\end{align*}
\color{black}

In practice the values $w_1, \ldots, w_n$ are never observed and the
estimators $\hat{\alpha}, \hat{\delta}$ are obtained with the EM-algorithm\footnote{Dempster, A.P., Laird, N.M., and Rubin, D.B. (1977) Maximum Likelihood
Estimation from Incomplete Data via the EM Algorithm, JRSS-B}.

With the EM algorithm $W_1, \ldots, W_n$ are considered to be missing and are 
estimated
by teir expected values $w_1^{(k)}, \ldots, w_n^{(k)}$ under the current
estimates $\alpha^{(k)}$ and $\delta^{(k)}$ (E-step). Then
$\ell(\alpha, \delta |y_1, \ldots, y_n, w_1^{(k)}, \ldots, w_n^{(k)})$ is 
maximized
(M-step), and the E- and M-steps are iterated until convergence.

More precisely, assuming $\alpha^{(k)}, \delta^{(k)}, w_1^{(k)}, \ldots,
w_n^{(k)}$ to be the values obtained at the $(k)-$th iteration, the iteration 
$(k+1)$
of the EM-algorithm requires the following two steps.
\begin{itemize}
  \item \textit{E-step.}  Find $w_i^{(k+1)}$ the conditional expectation of
$W_i$ given $y_i,  \alpha^{(k)}, \delta^{(k)}$ i.e. the expectation of $W_i$
with respect to the following probability mass function $P\{W_i|Y_i=y_i,
\alpha^{(k)}, \delta^{(k)}\}$.
  \item \textit{M-step.}  Find $\alpha^{(k+1)}, \delta^{(k+1)}$ as
$\argmax_{\alpha, \delta} \ell(\alpha, \delta |y_1, \ldots, y_n, w_1^{(k)},
\ldots, w_n^{(k)})$.
\end{itemize}

We would like you to provide the detailed description of the E-step, by 
answering the following
questions.
\begin{enumerate}
 \item[b)] Find the probability $P\{W_i|Y_i=y_i, \alpha^{(k)}, 
\delta^{(k)}\}$, for short
you can write it as $P\{W_i|y_i, \alpha, \delta\}$. In other words find
$P\{W_i=1|Y_i=0, \alpha^{(k)}, \delta^{(k)}\}$, $P\{W_i=0|Y_i=0,
\alpha^{(k)}, \delta^{(k)}\}$ and $P\{W_i=0|Y_i=y_i>0, \alpha^{(k)},
\delta^{(k)}\}$.

\noindent\textit{Hint: use the Bayes formula $P\{H|D\}={P\{D|H\} P\{H\}\over
P\{D\}}$.}\\  
 \\
\color{blue}
\begin{align*}
    P(W_i|y_i, \alpha, \delta) &= \frac{P(Y_i|w_i,\alpha, \delta)P(W_i|\alpha, \delta)}{P(Y_i=y_i|\alpha, \delta)}\\[1em]
    &= \frac{\left( (\exp\{-\lambda_i\}\lambda_i^{y_i}/y_i!)^{1-w_i}- \mathbbm{1}(y_i>0)\mathbbm{1}(w_i =1)\right) \pi_i^{w_i}(1-\pi_i)^{1-w_i}}{\left(\pi_i+(1-\pi_i)\exp(-\lambda_i)\right)^{\mathbbm{1}(y_i=0)}\left((1-\pi_i)\exp(-\lambda_i)\lambda_i^{y_i}/y_i!\right)^{\mathbbm{1}(y_i>0)}}\\[1em]
    &= \left\{
	\begin{array}{ll}
		\frac{\pi_i^{w_i}(1-\pi_i)^{1-w_i}\exp(-\lambda_i(1-w_i))}{\pi_i+(1-\pi_i)\exp(-\lambda_i)}  & \mbox{if } y_i=0; \\[1em]
		\frac{\pi_i^{w_i}(1-\pi_i)^{1-w_i}\left((\exp(-\lambda_i)\lambda_i^{y_i}/y_i!)^{1-w_i}-\mathbbm{1}(w_i=1)\right)}{(1-\pi_i)\exp(-\lambda_i)\lambda_i^{y_i}/y_i!} & \mbox{if } y_i>0.
	\end{array}
\right.
\end{align*}
In other words

\begin{align*}
    W_i =
\left\{
	\begin{array}{ll}
		1  & \mbox{with probability } \frac{\pi_i}{\pi_i+(1-\pi_i)\exp(-\lambda_i)}\ \mbox{if}\ y_i = 0 \\[1em]
		0  & \mbox{with probability } \frac{(1-\pi_i)\exp(-\lambda_i)}{\pi_i+(1-\pi_i)\exp(-\lambda_i)}\ \mbox{if}\ y_i = 0 \\[1em]
		0  & \mbox{with probability } 1\ \mbox{if}\ y_i > 0 \\
	\end{array}
\right.
\end{align*}
\color{black}

  \item[c)] Using $P\{W_i|Y_i, \alpha^{(k)}, \delta^{(k)}\}$ find
$w_i^{(k+1)}=E(W_i|Y_i=y_i,\alpha^{(k)}, \delta^{(k)})$.\\   \\
\color{blue}
From the previous point,
\begin{align*}
    w_i^{k+1}&=E\left(W_i|Y_i=y_i, \alpha^{(k)},\delta^{(k)}\right)\\[1em]
    &= \frac{\pi_i^{(k)}}{\pi_i^{(k)}+(1-\pi_i^{(k)})\exp(-\lambda_i^{(k)})}
\end{align*}
\color{black}

 \item[d)] How can we interpret the values of $w_1^{(k+1)}, \ldots,
w_n^{(k+1)}$ obtained at the last iteration of the algorithm ?\\  
\color{blue}
The values of $w_i^{(k+1)},...,w_n^{(k+1)}$ can show the probability of being from the spike of a zero.
\color{black}


 \end{enumerate}

\section{\underline{Exercise 3} }

{\color{blue} 
See the code for the solution.}

\begin{enumerate}
\item{Source the \texttt{truncpoisson.R} code available from the course website to create the \texttt{truncpoisson} GLM family object. Fit manually the two steps of the Poisson hurdle model to the data by using the \texttt{family=truncpoisson} argument in the \verb"glm" command. Compare your results to the output from the \texttt{hurdle} command of the \texttt{pscl} package.}
\item{Fit a simple Poisson GLM and a ZIP model. Among the three models, which provides the best fit?} 
\item{Do high insurance levels make the number of consultations increase?}
\end{enumerate}




\end{document}
